\section{Non-overlapping Intervals}
\label{sec:non-overlapping-intervals}

\problemstatement{We are given an array of intervals
$\textit{intervals}[i] = [\textit{start}_i, \textit{end}_i]$. We must find
the \textbf{minimum number of intervals to remove} so that the remaining
intervals are pairwise non-overlapping (two intervals that merely touch at
an endpoint, i.e.\ one ends exactly where another begins, are \emph{not}
considered overlapping).}

\begin{patternbox}
\emph{``Minimum number of intervals to remove so that the rest are
non-overlapping''} is logically the complement of
\emph{``maximum number of non-overlapping intervals you can keep''} -- and
the latter is exactly the activity selection problem from
Section~\ref{sec:n-meetings}. Whenever a problem asks for the
\textbf{minimum number to delete/remove} to enforce some property, check
whether the property is equivalent to a maximization you already know how
to solve; if so, the answer is simply
$n - (\text{the maximum achievable count})$. Recognizing this equivalence
saves you from re-deriving a new greedy strategy from scratch.
\end{patternbox}

\subsection*{Understanding the problem carefully}

Removing the fewest intervals to make the rest non-overlapping is the same
goal, viewed backward, as keeping the most intervals while staying
non-overlapping: every interval not removed must survive, and the survivors
must be mutually non-overlapping. So if we can compute
\[
  \textit{maxKept} = \text{largest number of pairwise non-overlapping
  intervals selectable},
\]
then the answer to this problem is immediately
\[
  n - \textit{maxKept}.
\]
And computing $\textit{maxKept}$ is precisely the activity selection problem
solved in Section~\ref{sec:n-meetings}: sort by finish time, and greedily
keep an interval whenever it does not conflict with the most recently kept
one.

\begin{ideabox}[Greedy Choice]
Sort intervals by end time, increasing. Walk through them in this order,
maintaining the end time of the last interval kept. Keep the current
interval (it does not need to be removed) if its start is at least the end
of the last kept interval (recall that touching endpoints are allowed, so
the comparison is $\ge$, not strictly $>$, unlike the room-booking version
in Section~\ref{sec:n-meetings} where touching endpoints \emph{do} conflict
because a physical room cannot instantly reset). Otherwise, this interval
must be removed, since it overlaps the last interval we chose to keep.
\end{ideabox}

\subsection*{Why this reduction and greedy choice are correct}

The correctness of greedily building $\textit{maxKept}$ by always preferring
the interval that finishes earliest among the remaining candidates is
exactly the exchange argument given in Section~\ref{sec:n-meetings}: the
interval with the globally earliest finish time can always be included in
some optimal non-overlapping selection, because including it can never
block more future options than including any later-finishing alternative
would. Once that first choice is fixed, the same argument applies
recursively to the remaining intervals that start after it ends. The only
change from Section~\ref{sec:n-meetings} is the boundary condition at
matching endpoints -- this problem explicitly states that touching is not
overlapping, so the compatibility check uses
$\textit{start}_i \ge \textit{lastEnd}$ rather than a strict inequality.

\subsection*{Algorithm}

\begin{enumerate}
  \item Sort $\textit{intervals}$ by end time, increasing.
  \item Initialize $\textit{lastEnd} \gets -\infty$ (or the end of the
        first interval after selecting it) and $\textit{kept} \gets 0$.
  \item For each interval $[\textit{s}, \textit{e}]$ in sorted order:
  \begin{enumerate}
    \item If $\textit{s} \ge \textit{lastEnd}$: keep this interval.
          $\textit{kept} \gets \textit{kept}+1$;
          $\textit{lastEnd} \gets \textit{e}$.
    \item Otherwise: this interval must be removed (it overlaps the last
          kept interval); do nothing further.
  \end{enumerate}
  \item Return $n - \textit{kept}$ as the minimum number of removals.
\end{enumerate}

\begin{algorithm}[H]
\caption{Minimum Removals for Non-overlapping Intervals}
\begin{algorithmic}[1]
\Procedure{EraseOverlapIntervals}{\textit{intervals}}
  \State sort \textit{intervals} by end time, increasing
  \State $\textit{lastEnd} \gets -\infty$, $\textit{kept} \gets 0$
  \For{each $[\textit{s}, \textit{e}]$ in sorted order}
    \If{$\textit{s} \ge \textit{lastEnd}$}
      \State $\textit{kept} \gets \textit{kept} + 1$
      \State $\textit{lastEnd} \gets \textit{e}$
    \EndIf
  \EndFor
  \State \Return $n - \textit{kept}$
\EndProcedure
\end{algorithmic}
\end{algorithm}

\subsection*{Dry run}

\dryrun{Take $\textit{intervals} = [[1,2],[2,3],[3,4],[1,3]]$.}

Sorted by end time: $[1,2],[2,3],[1,3],[3,4]$ (ties between intervals ending
at the same time can be broken arbitrarily; here $[2,3]$ and $[1,3]$ both
end at $3$).

\begin{itemize}
  \item $\textit{lastEnd}=-\infty$. Interval $[1,2]$: $1 \ge -\infty$.
        Keep. $\textit{kept}=1$, $\textit{lastEnd}=2$.
  \item Interval $[2,3]$: $2 \ge 2$. Keep (touching endpoints are allowed).
        $\textit{kept}=2$, $\textit{lastEnd}=3$.
  \item Interval $[1,3]$: $1 \ge 3$? No. Remove (do not update
        \textit{lastEnd} or \textit{kept}).
  \item Interval $[3,4]$: $3 \ge 3$. Keep. $\textit{kept}=3$,
        $\textit{lastEnd}=4$.
\end{itemize}

We kept $3$ out of $4$ intervals, so the minimum number of removals is
$4 - 3 = 1$ -- specifically, removing $[1,3]$ leaves
$[1,2],[2,3],[3,4]$, which are pairwise non-overlapping (touching only at
shared endpoints).

\subsection*{C++ implementation}

\begin{lstlisting}
#include <bits/stdc++.h>
using namespace std;

int eraseOverlapIntervals(vector<vector<int>>& intervals) {
    if (intervals.empty()) return 0;

    sort(intervals.begin(), intervals.end(),
         [](const vector<int>& a, const vector<int>& b) {
             return a[1] < b[1];
         });

    int n = intervals.size();
    int kept = 1;
    int lastEnd = intervals[0][1];

    for (int i = 1; i < n; i++) {
        if (intervals[i][0] >= lastEnd) {
            kept++;
            lastEnd = intervals[i][1];
        }
    }
    return n - kept;
}

int main() {
    vector<vector<int>> intervals = {{1,2},{2,3},{3,4},{1,3}};
    cout << "Minimum removals: "
         << eraseOverlapIntervals(intervals) << endl;
    return 0;
}
\end{lstlisting}

\begin{complexitybox}
\textbf{Time Complexity.} Sorting the $n$ intervals by end time costs
$O(n \log n)$. The subsequent single pass costs $O(n)$. The overall time
complexity is
\[
  O(n \log n).
\]
\textbf{Space Complexity.} Beyond the input array (sorted in place or with
$O(n)$ auxiliary space, depending on the sort implementation), we use only
a constant number of extra scalar variables, giving auxiliary space
complexity $O(1)$.
\end{complexitybox}

\begin{takeawaybox}
Many ``minimum removals/deletions to enforce property $P$'' problems are
secretly ``maximum subset satisfying property $P$'' problems wearing a
disguise: the answer is $n$ minus the size of the largest valid subset.
Whenever you see a removal-minimization phrasing, check whether flipping it
to a keep-maximization phrasing turns it into a pattern you already
recognize -- here, it turns directly into the activity selection problem
from Section~\ref{sec:n-meetings}.
\end{takeawaybox}
